\section{Riproducibilità computazionale e provenienza}
\label{app:repro}

\subsection{Provenienza della campagna finale di twin experiments}

Il target sintetico finale è stato riutilizzato senza rigenerazione e verificato tramite SHA-256 \texttt{1bca9e29\allowbreak{}c60d77c1\allowbreak{}450fd2fc\allowbreak{}ea845d23\allowbreak{}20c9e449\allowbreak{}45dd7d3c\allowbreak{}cda52aa3\allowbreak{}3ef261af}. La discretizzazione target è $(N,M)=(32,160)$, quella inversa principale $(N,M)=(16,80)$ e il numero di osservazioni è 180. La baseline senza rumore riutilizza tre run con seed CPSO 20261017, 20261018 e 20261019. La campagna rumorosa usa il seed CPSO 20261017 e i seed dati $(2026110501,\ldots,2026110505)$ per $\eta=0.05$ e $(2026111001,\ldots,2026111005)$ per $\eta=0.10$. Il protocollo pianifica 600 valutazioni FEM per run; gli artefatti JSON registrano la storia completa, il costo migliore, la condizione di arresto e i controlli di integrità.

Il test di integrità della campagna finale non sostituisce una validazione indipendente del modello fisico né una verifica di convergenza dell'ottimizzatore. Prima della sottomissione, l'archivio pubblico versionato dovrà contenere i sorgenti effettivamente usati, le configurazioni, le dipendenze, gli artefatti JSON/CSV e i log necessari per riprodurre i numeri inclusi nel manoscritto.

\subsection{Ambiente software e suite di test}

Il report di provenienza archiviato per il commit \texttt{2c10b387833dd81fbae3654597bdcad823cacfeb} registra 37 test superati in 57.66 s nell'ambiente dichiarato: macOS 26.5.1 su arm64, Python 3.11.17 e deal.II indicato nel file di ambiente. Questo conteggio va associato alla specifica snapshot documentata dal report; prima della sottomissione, la suite deve essere rieseguita sul commit esatto destinato al rilascio. Il file \texttt{environment\_provenance.json}, il lock Conda, i report JSON, le matrici Jacobiane, le tabelle CSV, le mesh e i log originali sono inclusi nell'archivio sorgente consegnato con il manoscritto.

I test elencati nel report comprendono formule e limiti dell'operatore L1, compatibilità iniziale, identità energetica e massa discreta, MMS accoppiato, non-identificabilità a flusso nullo, controlli della pipeline anatomica, integrazione di mesh quadrilatere, retrocompatibilità MMS, norma vettoriale $L^2$ ed entry point CLI. Il passaggio dei test supporta la riproducibilità del comportamento verificato nei casi coperti, ma non dimostra da solo correttezza fisica generale o validità clinica.

\subsection{Provenienza del caso anatomico}

I report della geometria contengono SHA-256 e percorsi degli input MRI effettivamente usati. L'identità dei file e la griglia NIfTI vengono quindi tracciate nel bundle; la provenienza ufficiale del dataset deve essere attestata tramite la fonte BraTS e le condizioni di accesso, non dal solo hash. Il soggetto è un singolo caso del benchmark. Le forzanti, i campi temporali e i dati per i twin experiments sono sintetici.
